%%%%%%%%%%%%%%%%%%%%%%%%%%%%%%%%%%%%%%%%%%%%%%%%%%%%%%%%%%%%%%%%%%%%%%%%%%%%%%%%
% Chapter 1
%%%%%%%%%%%%%%%%%%%%%%%%%%%%%%%%%%%%%%%%%%%%%%%%%%%%%%%%%%%%%%%%%%%%%%%%%%%%%%%%

\chapter{Why Covariant Stochastic Mechanics?\\
Historical, Physical and Mathematical Motivation}

\epigraph{%
``The important thing is not to stop questioning.''%
}{Albert Einstein}

\section{Introduction}

Quantum mechanics and the theory of relativity are unquestionably among the
greatest intellectual achievements of twentieth-century physics. Their
predictions have been confirmed to extraordinary precision over an enormous
range of scales, from atomic spectroscopy to high-energy particle collisions
and from gravitational lensing to the dynamics of compact astrophysical
objects. Nevertheless, despite these empirical successes, their conceptual
unification remains incomplete.

The difficulty is not merely technical. Quantum field theory provides a
remarkably successful computational framework combining quantum principles
with Lorentz invariance, yet many of the foundational questions raised during
the first decades of quantum mechanics remain essentially unchanged.
What is the physical meaning of the wave function? Does quantum mechanics
describe objective physical reality or merely probabilities for observations?
Can one construct a relativistic theory possessing well-defined particle
trajectories? What is the precise role of time in quantum theory? Why do
probabilities arise from complex probability amplitudes rather than ordinary
stochastic processes?

These questions acquire additional significance in relativistic physics.
Whereas nonrelativistic quantum mechanics possesses a preferred Newtonian
time parameter, relativistic space-time treats temporal and spatial
coordinates on nearly equal footing. Consequently, many constructions that are
mathematically natural in the Schrödinger theory become problematic after the
imposition of Lorentz covariance. Examples include localization of particles,
the definition of probability densities, and perhaps most importantly the
definition of stochastic processes themselves.

The present monograph is motivated by the hypothesis that these apparently
independent difficulties may have a common origin. We shall investigate the
possibility that relativistic quantum mechanics emerges from an underlying
covariant stochastic dynamics formulated not with respect to coordinate time
but with respect to an invariant evolution parameter. Such a formulation is
naturally suggested by the manifestly covariant relativistic dynamics
developed independently by Stueckelberg, Horwitz, Piron and their
collaborators. Within that framework, evolution is parametrized by a universal
invariant parameter, usually denoted by $\tau$, while space-time coordinates
become dynamical variables on an equal footing.

The central question investigated throughout this monograph may therefore be
stated as follows:

\begin{center}
\fbox{
\begin{minipage}{0.92\textwidth}
\begin{center}

\emph{Can relativistic quantum mechanics be understood as the statistical
mechanics of covariant stochastic worldlines evolving in invariant time?}

\end{center}
\end{minipage}}
\end{center}

This question immediately separates the present work from most existing
approaches to stochastic mechanics. Traditional stochastic formulations,
beginning with the pioneering works of Fényes and Nelson, employ the ordinary
Newtonian time variable as the evolution parameter of the stochastic process.
Although remarkably successful in reproducing the Schrödinger equation, such
formulations encounter severe conceptual difficulties when generalized to
special relativity. The principal obstacle is that coordinate time is not a
Lorentz-invariant quantity and therefore cannot naturally serve as the
parameter of a covariant diffusion process.

Manifestly covariant relativistic dynamics suggests a possible resolution.
Instead of employing coordinate time, one introduces an invariant evolution
parameter that labels successive events along the dynamical history of the
system. The resulting formalism possesses a canonical Hamiltonian structure
closely analogous to ordinary classical mechanics while preserving manifest
Lorentz covariance. From the viewpoint of stochastic analysis, this invariant
parameter appears to be a natural candidate for the role traditionally played
by Newtonian time in diffusion theory.

The purpose of the present monograph is not to advocate a particular
interpretation of quantum mechanics from the outset. Rather, our objective is
to examine systematically whether the mathematical structures of stochastic
analysis and manifestly covariant Hamiltonian dynamics can be unified into a
single coherent framework. Such a framework should satisfy several demanding
requirements simultaneously.

First, it should reproduce ordinary nonrelativistic stochastic mechanics in
the appropriate limit.

Second, it should respect manifest Lorentz covariance at every stage of the
construction.

Third, it should admit a consistent treatment of interacting many-particle
systems.

Fourth, it should provide a natural geometric setting for internal degrees of
freedom such as spin.

Finally, it should reduce to the experimentally successful predictions of
relativistic quantum theory whenever the latter has been tested.

Whether such a program can ultimately be completed remains an open question.
Nevertheless, enough mathematical developments have accumulated during the
past seventy years to justify a comprehensive re-examination of the problem.
The present volume attempts to provide both a detailed review of these
developments and the foundations upon which a more general theory may
eventually be constructed.
%%%%%%%%%%%%%%%%%%%%%%%%%%%%%%%%%%%%%%%%%%%%%%%%%%%%%%%%%%%%%%%%%%%%%%%%%%%%%%
\section{Why Another Foundation of Quantum Mechanics?}
%%%%%%%%%%%%%%%%%%%%%%%%%%%%%%%%%%%%%%%%%%%%%%%%%%%%%%%%%%%%%%%%%%%%%%%%%%%%%%

The extraordinary empirical success of quantum mechanics has often been
interpreted as evidence that its conceptual foundations are essentially
complete. From this perspective, the remaining open problems concern only
technical extensions such as quantum gravity, strongly correlated systems, or
the mathematical properties of quantum field theory. Such an assessment,
however, overlooks an important distinction between predictive success and
conceptual completeness.

A physical theory may provide an extraordinarily accurate algorithm for the
calculation of observable quantities while nevertheless leaving unanswered
questions concerning the ontology of the systems it describes. The history of
physics contains numerous examples. Ptolemaic astronomy successfully predicted
planetary positions for many centuries without providing a correct dynamical
description of the Solar System. Classical thermodynamics achieved remarkable
accuracy long before the atomic constitution of matter became generally
accepted. Likewise, Maxwell's electromagnetic theory preceded the microscopic
understanding of electric charge by several decades.

The existence of unresolved conceptual questions should therefore not be
interpreted as evidence against the empirical validity of quantum mechanics.
Rather, it motivates continued investigation into whether the present
formalism admits a deeper dynamical explanation.

Among the questions that have persisted since the birth of quantum mechanics,
the following remain particularly significant.

\begin{enumerate}

\item
Why should the fundamental object describing an individual physical system be
a complex-valued wave function rather than ordinary classical observables?

\item
Why does the Born probability density arise as the squared modulus of the wave
function?

\item
Can quantum probabilities be understood as emerging from an underlying
objective stochastic dynamics?

\item
Does a relativistic particle possess a well-defined space-time trajectory
between measurements?

\item
What is the physical meaning of quantum nonlocal correlations?

\item
What is the status of time within quantum theory?

\end{enumerate}

It is important to emphasize that these questions are logically independent of
the experimental confirmation of quantum mechanics. They concern the
interpretation and mathematical structure of the theory rather than its
phenomenological accuracy.

The present monograph adopts a deliberately conservative attitude toward these
questions. Rather than modifying the experimentally successful predictions of
quantum theory, we seek to determine whether they may emerge from a more
primitive dynamical framework possessing the following properties:

\begin{enumerate}

\item manifest Lorentz covariance,

\item Hamiltonian dynamics,

\item stochastic evolution,

\item objective space-time histories,

\item compatibility with known quantum phenomena.

\end{enumerate}

The possibility that such a framework exists cannot be dismissed a priori.
Indeed, each of these ingredients already appears separately within modern
mathematical physics. The principal difficulty has been to combine them into a
single coherent theory.
%%%%%%%%%%%%%%%%%%%%%%%%%%%%%%%%%%%%%%%%%%%%%%%%%%%%%%%%%%%%%%%%%%%%%%%%%%%%%%%%
\section{What Exactly Is the Problem?}
%%%%%%%%%%%%%%%%%%%%%%%%%%%%%%%%%%%%%%%%%%%%%%%%%%%%%%%%%%%%%%%%%%%%%%%%%%%%%%%%

The search for a relativistic stochastic formulation of quantum mechanics is
often presented as an attempt to reconcile two successful theories:
relativity and quantum mechanics. While this statement is historically
accurate, it conceals a more fundamental issue. Neither theory is internally
problematic within its own domain of applicability. Rather, the difficulty
arises because they embody profoundly different conceptions of physical
description.

Classical mechanics is formulated in terms of individual systems. The
fundamental dynamical object is a trajectory in configuration space or phase
space, and the equations of motion determine how this trajectory evolves in
time from specified initial conditions. Probabilities enter only through
ignorance of the precise initial state or through interaction with an
environment.

Quantum mechanics adopts a different standpoint. The primary object is not a
trajectory but a vector in a Hilbert space, or more generally a density
operator. Observable quantities are represented by self-adjoint operators,
while probabilities are intrinsic components of the formalism rather than
expressions of incomplete information. Although this framework has proved
extraordinarily successful, it leaves open the question of whether the
Hilbert-space description is fundamental or emergent.

Relativity introduces an additional conceptual shift. In Newtonian mechanics,
time is universal and absolute. Every observer agrees on the temporal order of
events, and every dynamical process evolves with respect to the same parameter.
Special relativity replaces this universal temporal ordering by a four-
dimensional space-time in which temporal and spatial coordinates transform
under the Lorentz group. The absence of an observer-independent coordinate
time has profound consequences for any dynamical theory based upon temporal
evolution.

The difficulty becomes particularly apparent in stochastic mechanics. A
stochastic process is traditionally defined with respect to an increasing
parameter possessing a natural ordering. For diffusion processes this parameter
is interpreted as time. The mathematical construction depends fundamentally
upon this ordering through filtrations, conditional expectations and Markov
semigroups.

In nonrelativistic physics the identification of the stochastic parameter with
Newtonian time presents no conceptual difficulty. In relativistic physics,
however, coordinate time depends upon the observer. Consequently, the very
definition of a Lorentz-covariant stochastic process becomes nontrivial.

This observation is frequently overlooked. The principal obstacle is not the
construction of relativistic diffusion coefficients nor the covariance of the
resulting differential equations. Rather, it concerns the identification of
the parameter with respect to which stochastic evolution itself is defined.

Manifestly covariant relativistic dynamics provides a remarkably elegant
alternative. Instead of employing coordinate time as the evolution parameter,
one introduces an invariant world parameter, conventionally denoted by
$\tau$. Space-time coordinates then become ordinary dynamical variables,
placing temporal and spatial coordinates on equal footing.

From the perspective of stochastic analysis this change is profound. The
parameter $\tau$ possesses precisely the mathematical properties required of
the evolution parameter of a Markov process: it is monotonically increasing
along every dynamical history, invariant under Lorentz transformations, and
independent of any particular inertial observer.

The existence of such a parameter suggests that the apparent incompatibility
between stochastic dynamics and relativity may not be fundamental. Instead,
it may arise from the historically inherited assumption that stochastic
processes must necessarily evolve with respect to coordinate time.

The central thesis explored throughout this monograph is therefore not that
quantum mechanics should be modified, but rather that the kinematical
framework within which stochastic dynamics is formulated may require
reconsideration.

If invariant evolution replaces coordinate time as the fundamental dynamical
parameter, the principal conceptual obstacle that has impeded relativistic
stochastic mechanics for more than half a century may disappear. Whether this
expectation can be realized consistently remains one of the principal
questions investigated in the chapters that follow.


